\cvskill{Italiano [Italian]}{5}
My {\em mother tongue}. In addition to the normative form, I know my regional dialect as well.

\divider

\cvskill{English}{4}
Model/variety: modern/contemporary RP. \\
CEFR C1--C2 proficiency level.
%{\em B2 Certificate}. Center for Applied Linguistics (CLA) of UNIPG, Perugia (Italy), July 2017

\divider

\cvskill{Esperanto [Esperanto]}{4}
{\em CEFR C1}.

\divider

\cvskill{Deutsch [German]}{3}
Model/variety: standard High German. \\
{\em CEFR B1} language-proficiency certificate, {\it Centre for Applied Linguistics} (CLA) of UNIPG, Perugia (Italy), July 2017.

\divider

\cvskill{Castellano [Spanish]}{2}
{\em CEFR A2} when speaking/writing, {\em B1--B2} when listening/reading.

\divider

\cvskill{日本語 [Japanese]}{1}
{\em CEFR A1--A2}. I have got a smattering of the language; I can hold very basic and short conversations around a few topics.

\divider

\cvskill{{\hindifont हिन्दी } [Hindi]}{1}
{\em CEFR A1}. Basic phrases and grammar.


\newpage

\cvsection{Software Projects}

\cvevent{Algorithms and Data Structures for Efficient Ride-Sharing Platforms}{Master's thesis research fellow}{Apr 2019 -- Feb 2020}{Pisa, Italy}

Implementation of a novel locality filtering technique which can speed up matching computations for urban mobility problems. The proposed algorithm and data structure solve the ride-sharing problem by significantly improving both time and space complexities. We developed this project in collaboration with researchers at {\em CNR} (Pisa, Italy) and {\em MIT Senseable City Lab} (Boston, Massachusetts, USA). This work was accepted for publication on {\em IEEE Transactions on Intelligent Transportation Systems}.

\divider

\cvevent{Parallel Versions of the Particle-Swarm Optimisation Method}{Parallel and distributed programming course}{Fall 2019}{Pisa, Italy}

Study for a map-reduce parallelisation of the PSO procedure for the localization of a minimum within a 2-D plane. The parallel versions have been realized using: (1) \texttt{c++} threads, as well as (2) the FastFlow parallel framework. We executed our experiments on a {\em Intel Xeon Phi} machine (64 cores, 4-way hyper-threading).

\divider

\cvevent{SPADE -- Sensing, Processing and Analysing Data in an Embedded IoT Device}{Wireless networks of embedded systems course}{Spring 2019}{Pisa, Italy}

Design and deployment of a machine-learning powered system for the predictive maintenance of machines based on a Beagleboard green and an accelerometer.

\divider

\cvevent{Implementation of a Distributed Graph-Clustering Algorithm}{Bachelor's thesis}{Fall 2017}{Perugia, Italy}

Distributed implementation of the LPAm clustering algorithm on top of the Apache Giraph platform. The programming paradigm of the Giraph library is ``Think Like a Vertex'' (TLAV), which aims to exploit better locality in data access and communications.


\begin{comment}
\cvsection{s}
\begin{itemize}
    \item 
    Language certificates (English and German)
    
    \item
    ECDL certificate
    
    \item
    EPLASS and Flyeralarm internship certificates
    
    \item
    BSc\ and MSc\ courses transcripts
    
    \item
    master's thesis
    
\end{itemize}
\end{comment}


% Adapted from @Jake's answer from http://tex.stackexchange.com/a/82729/226
% \wheelchart{outer radius}{inner radius}{
% comma-separated list of value/text width/color/detail}
%\wheelchart{1.5cm}{0.5cm}{%
%  6/8em/accent!30/{Sleep,\\beautiful sleep},
%  3/8em/accent!40/Hopeful novelist by night,
%  8/8em/accent!60/Daytime job,
%  2/10em/accent/Sports and relaxation,
%  5/6em/accent!20/Spending time with family
%}

%% Yeah I didn't spend too much time making all the
%% spacing consistent... sorry. Use \smallskip, \medskip,
%% \bigskip, \vpsace etc to make adjustments.


%% Switch to the right column. This will now automatically move to the second
%% page if the content is too long.
%\switchcolumn



%%%%%%%%%%%%%%%%%%%%%%%%%%%%%%%%%%%%%%%%%%%%%%%%%%

    \euro{}1500 – Awarded the \textbf{``Best Graduate Award'' (Master’s in Computer Science \& Networking)} for achieving the highest graduation grade, outstanding exam performance (GPA), and the shortest time to graduation. This award was reserved for students from the 2017/2018 cohort who graduated in the academic year 2018/2019.